\documentclass[PMI,VKR]{HSEUniversity}
% Возможные опции: KR или VKR; PI или BI

\title{Моделирование хаотических динамических систем с помощью NeuralODE}
\author{Мяков Тимофей Ильич}
\Year{2026}
\Abstract{
В работе исследуется применение моделей семейства Neural Ordinary Differential Equations к задаче восстановления и прогнозирования нелинейных динамических систем по временным рядам. 
Особое внимание уделяется хаотическим и квазипериодическим режимам, для которых малые ошибки в оценке векторного поля могут приводить к быстрому расхождению траекторий. 
В качестве тестовых систем рассматриваются система Рёсслера, модель Хиндмарша--Роуза и квазипериодический осциллятор, демонстрирующие соответственно хаотическую, берстовую и квазипериодическую динамику.

В работе сравниваются классическая NeuralODE, Augmented NeuralODE, ControlSynth NeuralODE, а также предложенная архитектура FNODE. 
FNODE объединяет расширение фазового пространства, диссипативный линейный член в векторном поле, компактное нелинейное MLP-ядро и специальный протокол обучения, включающий curriculum learning, teacher forcing и спектральную функцию потерь. 
Обучающие данные формируются численным интегрированием эталонных систем методом Рунге--Кутты четвёртого порядка, после чего модели обучаются на окнах траекторий и оцениваются на отложенных временных фрагментах.

Экспериментальная оценка проводится по метрикам MSE, MAE, дивергенции Кульбака--Лейблера, времени валидности прогноза и специализированным характеристикам берстовой динамики. 
Результаты показывают, что предложенная FNODE демонстрирует наилучшее или близкое к наилучшему качество на всех рассмотренных системах. 
На системах Рёсслера и квазипериодического осциллятора FNODE существенно превосходит базовые подходы по точности прогноза и воспроизведению распределения состояний, а на системе Хиндмарша--Роуза обеспечивает конкурентоспособную ошибку и более точное восстановление берстового режима. 

Дополнительные эксперименты с зашумлёнными данными и различным объёмом обучающей выборки подтверждают устойчивость предложенного подхода и его применимость к моделированию сложной нелинейной динамики.
}

\bibliography{library.bib}

\begin{document}

\maketitle


\chapter{Введение}

Моделирование динамических систем - важная современная задача для прикладных областей науки. 
Многие реальные процессы в физике, биологии,
нейронауках и инженерии описываются системами дифференциальных
уравнений, демонстрирующими сложное поведение: устойчивые колебания,
бифуркации, квазипериодичность и хаос \cite{strogatz2018nonlinear, ott2002chaos}.
Особую трудность представляет прогнозирование хаотических систем,
поскольку даже малые изменения в начальных условиях или в оценке
векторного поля приводят к быстрому расхождению траекторий
\cite{kantz2004nonlinear}.

Классический подход к моделированию динамических систем предполагает
задание уравнений движения на основе физических законов или
экспертного знания. Однако в реальных приложениях точная форма
уравнений часто неизвестна, а доступны только наблюдения временного
ряда. В таких условиях возникает задача восстановления динамики по
данным. Один из современных подходов к этой задаче связан с
использованием нейронных сетей для аппроксимации правой части
дифференциального уравнения. В частности, модели Neural Ordinary
Differential Equations позволяют интерпретировать нейронную сеть как
непрерывную динамическую систему и обучать параметризованное векторное
поле непосредственно по наблюдаемым траекториям
\cite{chen2019neuralordinarydifferentialequations}.

Цель данной работы состоит в исследовании возможностей семейства NeuralODE
архитектур при моделировании хаотических и квазипериодических
динамических систем. Основное внимание уделяется не только
точности краткосрочного прогноза, но и способности моделей
воспроизводить долгосрочную статистическую структуру аттрактора,
сохранять частотные характеристики и демонстрировать устойчивость к
шуму в данных.

В качестве тестовых примеров рассматриваются система Рёсслера,
система Хиндмарша-Роуза и модель генератора квазипериодических колебаний. Эти системы
демонстрируют различные типы динамического поведения и потому позволяют
проверить модели в нескольких режимах: хаотическом, берстовом и
квазипериодическом. Такой выбор систем делает экспериментальную часть
работы более репрезентативной и позволяет оценить сильные и слабые
стороны рассматриваемых архитектур.


\chapter{Обзор литературы}

\section{Методы Neural Ordinary Differential Equations}

\subsection{NeuralODE}

Модели семейства \textit{Neural Ordinary Differential Equations},
впервые предложенные в работе Chen et al.~\cite{chen2019neuralordinarydifferentialequations},
предлагают непрерывную интерпретацию глубокой нейронной сети.
Связь между residual-сетями и численными схемами интегрирования
динамических систем также обсуждалась в работах
\cite{haber2017stable, ruthotto2020deep}. В классических архитектурах
преобразование скрытого состояния задаётся последовательностью
дискретных слоёв. Например, в residual-сетях переход между соседними
слоями имеет вид
\[
h_{k+1}=h_k+f(h_k,\theta_k).
\]
Такой шаг можно рассматривать как численную аппроксимацию динамической
системы методом Эйлера. Если уменьшать шаг дискретизации и одновременно
увеличивать число слоёв, то в пределе возникает непрерывная по
``глубине'' модель, описываемая обыкновенным дифференциальным
уравнением
\[
\frac{dh(t)}{dt}=f(h(t),t,\theta).
\]
Здесь функция $f$ параметризуется нейронной сетью и определяет
векторное поле в пространстве скрытых состояний. По начальному значению
$h(t_0)$ решение уравнения в момент времени $t_1$ вычисляется как
\[
h(t_1)=\mathrm{ODESolve}(h(t_0),f,t_0,t_1,\theta),
\]
где $\mathrm{ODESolve}$ обозначает внешний численный решатель ОДУ.

Таким образом, вместо явного задания конечного набора слоёв модель
задаёт непрерывную траекторию преобразования признаков. Это позволяет
трактовать глубину сети не как фиксированное число блоков, а как
интервал интегрирования динамической системы.

\begin{figure}[H]
  \centering
    \includegraphics[width=1\textwidth]{./img/theory/node_func.png}
    \caption{Схема обучения семества моделей NeuralODE}
    \label{fig:node_func}
\end{figure}

\subsection{Augmented NeuralODE}

Основная идея \textit{Augmented NeuralODE} была предложена в работе
Dupont et al.~\cite{dupont2019augmented}. Она заключается в дополнении
скрытого состояния координатами, которые не обязательно имеют
непосредственную интерпретацию в исходном пространстве признаков, но
позволяют динамической системе реализовывать более широкий класс
преобразований.

Исходное состояние \(h(t)\in\mathbb{R}^d\) расширяется дополнительным
вектором \(a(t)\in\mathbb{R}^p\), после чего динамика задаётся уже в
пространстве большей размерности:
\[
z(t)=
\begin{bmatrix} 
    h(t)\\ a(t)
\end{bmatrix}
\in\mathbb{R}^{d+p}.
\]
Обычно начальное значение дополнительных координат выбирается равным
нулю:
\[
z(t_0)=
\begin{bmatrix}
    h(t_0) \\ 0
\end{bmatrix}.
\]
Далее \(z(t)\) используется как аналогично простой NeuralODE

После интегрирования для дальнейшего предсказания может использоваться
всё расширенное состояние \(z(t_1)\) либо только его часть,
соответствующая исходным признакам:
\[
h(t_1)=\pi_h(z(t_1)),
\]
где \(\pi_h\) обозначает проекцию на исходные координаты.

Благодаря этому модель может реализовывать преобразования, недоступные
классической NODE без аугментации. С геометрической точки зрения
Augmented NODE предоставляет динамической системе дополнительные
степени свободы: данные могут временно ``выходить'' в дополнительные
координаты, изменять своё взаимное расположение, а затем
проецироваться обратно.

\subsection{Control-Synth NeuralODE}

\textit{ControlSynth Neural Ordinary Differential Equations}\cite{mei2024controlsynthneuralodesmodeling} были предложены как расширение классических NeuralODE для моделирования сложных динамических систем с гарантируемой сходимостью. 
В ControlSynth NODE динамика строится как более структурированная система, включающая как основное нелинейное векторное поле, так и дополнительный управляющий член. 
Такая конструкция заимствует идеи из теории управления и позволяет не только аппроксимировать наблюдаемую динамику, но и накладывать на неё дополнительные свойства устойчивости и сходимости.


В ControlSynth NODE задаётся как управляемая динамическая система
\[
\frac{dx(t)}{dt}=F(x(t),t,\theta,\phi),
\]
где правая часть имеет составную структуру:
\[
F(x(t),t,\theta,\phi)
=
f_{\theta}(x(t),t)+g_{\phi}(x(t),t).
\]
Здесь \(f_{\theta}\) отвечает за базовое описание нелинейной динамики, а \(g_{\phi}\) играет роль дополнительного управляющего или корректирующего воздействия. За счёт такого разложения модель получает возможность отдельно учитывать основное течение процесса и дополнительные факторы, влияющие на его эволюцию.

Решение системы, как и в классических NeuralODE, вычисляется с помощью внешнего численного решателя:
\[
x(t_1)=\mathrm{ODESolve}(x(t_0),F,t_0,t_1,\theta,\phi).
\]
Таким образом, выход модели определяется не последовательностью дискретных слоёв, а непрерывной траекторией состояния, полученной при интегрировании управляемого дифференциального уравнения.

Ключевая особенность ControlSynth NODE состоит в том, что её структура позволяет анализировать сходимость модели с помощью условий из теории управления. 
В частности, несмотря на нелинейный характер правой части, сходимость может быть обеспечена за счёт проверяемых линейных неравенств. Это отличает ControlSynth NODE от обычных NODE, где устойчивость и поведение траекторий часто являются следствием только процесса обучения и не гарантируются явно.

\section{Динамические системы}

Далее представлены динамические системы, на которых будут построены датасеты для обучения моделей. 
Стоит заметить, что для этих систем существует описание в виде системы ОДУ, поэтому мы можем получить достаточно качественные наборы данных.
В реальных задачах, скорее всего, у нас будет лишь временной ряд, с которым предстоит работать.

\subsection{Система Рёсслера}

Система Ресслера является одной из эталонных моделей нелинейной динамики, в которой наблюдается хаотическое поведение. 
Систем была предложена Отто Ресслером в 1976 \cite{rossler1976equation} как динамическая система, описывающая изменение концентрации веществ в химических реакциях, 
и представляет собой систему трех обыкновенных дифференциальных уравнений. 
Система Ресслера является классическим примером динамической системы, в которой наблюдается развитие хаоса через каскад бифуркаций удвоения периода.

Модель может быть описана системой ОДУ~\eqref{eq:rossler_ode}.

\begin{figure}[h]
    \centering

    \begin{minipage}[c]{0.45\textwidth}
        \centering
        \begin{equation}
        \label{eq:rossler_ode}
        \begin{cases}
            \dfrac{dx}{dt} = -y - z \\[10pt]
            \dfrac{dy}{dt} = x + ay \\[10pt]
            \dfrac{dz}{dt} = b + z(x - c)
        \end{cases}
        \end{equation}
        
        \vspace{-0.5em}
        \[
            a = 0.2, \quad b = 0.2, \quad c = 5.7.
        \]
    \end{minipage}
    \hfill
    \begin{minipage}[c]{0.45\textwidth}
        \centering
        \includegraphics[width=\textwidth]{./img/theory/rossler_phase_portrait.png}
        \caption{Фазовый портрет системы Рёсслера}
    \end{minipage}

\end{figure}

\subsection{Система Хиндмарша-Роуза}

Система Хиндмарша-Роуза описывает электрическую активность нейрона и
является одной из классических низкоразмерных моделей для изучения
нейронной динамики \cite{hindmarsh1984model, rose1985model}. Она была
предложена для качественного описания спайковой и берстовой активности,
наблюдаемой в экспериментах с отдельными нейронами.

Система может быть описана системой ОДУ~\eqref{eq:hr_ode}. Здесь
переменная $x(t)$ интерпретируется как мембранный потенциал нейрона,
$y(t)$ описывает быструю восстановительную переменную, а $z(t)$ —
медленную адаптационную переменную, отвечающую за переходы между
активными и покоящимися фазами берста
\cite{hindmarsh1984model, izhikevich2007dynamical}.

\begin{figure}[h]
    \centering

    \begin{minipage}[c]{0.45\textwidth}
        \begin{equation}
        \label{eq:hr_ode}
        \begin{cases}
            \dfrac{dx}{dt} = y - ax^3 + bx^2 - z + I \\[10pt]
            \dfrac{dy}{dt} = c - dx^2 - y \\[10pt]
            \dfrac{dz}{dt} = r\left[s(x - x_0) - z\right]
        \end{cases}
        \end{equation}
        \vspace{-0.5em}
        \[
            a=1.0, \quad b=3.0, \quad c=1.0, \quad d=5.0
        \]
        \[
            I=3.25, \quad r=0.006
        \]
        \[
            s=4.0,\quad  x_0=-1.6
        \]
    \end{minipage}
    \hfill
    \begin{minipage}[c]{0.45\textwidth}
        \includegraphics[scale=0.3]{./img/theory/hr_phase_portrait.png}
        \caption{Фазовый портрет системы Хиндмарша-Роуза (численный метод — RK4)}
    \end{minipage}

\end{figure}

\subsection{Квазипериодический осциллятор}

Квазипериодическая динамика является одним из типичных режимов
нелинейных автоколебательных систем и часто связана с движением на
инвариантном торе в фазовом пространстве
\cite{guckenheimer1983nonlinear, strogatz2018nonlinear}.

Рассматриваемая система представляет модель радиофизического генератора типа ван дер Поля с трехмерным фазовым пространством, 
в которой можно наблюдать последовательность бифуркаций, приводящих к формированию автономных квазипериодических колебаний, 
в фазовом пространстве которым будет отвечать тор.
 Параметр \(\lambda\)
отвечает за уровень возбуждения автоколебаний, \(\omega_0\) задаёт
собственную частоту линейной части осциллятора, \(\beta\) определяет
силу нелинейного насыщения, а \(\kappa\) — интенсивность обратной связи
через переменную \(z\). При выбранных наборах параметров система
демонстрирует сложные устойчивые колебательные режимы, в том числе
квазипериодическое поведение. Модель может быть описана системой ОДУ~\eqref{eq:qp_ode}

\begin{figure}[h]
    \centering

    \begin{minipage}[c]{0.45\textwidth}
        \begin{equation}
        \label{eq:qp_ode}
        \begin{cases}
            \dfrac{dx}{dt} = y \\[10pt]
            \dfrac{dy}{dt} = (\lambda + z + x^2 - \beta x^4)\,y - \omega_0^2 x \\[10pt]
            \dfrac{dz}{dt} = -z - \kappa y^2
        \end{cases}
        \end{equation}
        \vspace{-0.5em}
        \[
            \lambda=0.5, \quad \beta=1/18,
        \]
        \[
            \kappa=0.02, \quad \omega_0=5.1
        \]
        
    \end{minipage}
    \hfill
    \begin{minipage}[c]{0.45\textwidth}
        \includegraphics[scale=0.3]{./img/theory/qp_phase_portrait.png}
        \caption{Фазовый портрет квазипериодического генератора}
    \end{minipage}

\end{figure}

Данные системы используются как тестовые динамические системы при
исследовании методов восстановления и прогнозирования нелинейной
динамики. Они позволяют проверять, насколько хорошо модель способна
воспроизводить не только амплитуду и частоту колебаний, но и более
сложную структуру фазовой траектории, характерную для
хаотических/квазипериодических процессов.


\chapter{Методология}

\section{Подготовка данных для обучения}

Для всех динамических систем моделировались ground truth траектории с
помощью численного метода RK4. 

Метод Рунге-Кутты 4-го порядка (RK4) является одним из наиболее популярных явных одношаговых методов численного интегрирования. Он обеспечивает четвертый порядок точности по шагу $h$.


Алгоритм подготовки данных для обучения:

\begin{algorithm}[H]
\caption{Сбор траектории методом RK4}
\label{alg:rk4_trajectory}
\begin{algorithmic}[1]
\Require модель правой части $f(t, \mathbf{y})$, начальное состояние $\mathbf{y}_0$, 
начальное время $t_0$, конечное время $T$, шаг интегрирования $h$
\Ensure траектория $\{\mathbf{y}_n\}_{n=0}^{N}$ в узлах сетки $\{t_n\}_{n=0}^{N}$

\State $N \gets \lceil (T - t_0) / h \rceil$
\State Инициализировать массив $\mathbf{Y} \leftarrow [\mathbf{y}_0]$
\State $t \gets t_0$
\State $\mathbf{y} \gets \mathbf{y}_0$
\State
\For{$n = 0, \dots, N-1$}
    \State $\mathbf{k}_1 \leftarrow f(t, \mathbf{y})$
    \State $\mathbf{k}_2 \leftarrow f(t + h/2, \mathbf{y} + h \cdot \mathbf{k}_1 / 2)$
    \State $\mathbf{k}_3 \leftarrow f(t + h/2, \mathbf{y} + h \cdot \mathbf{k}_2 / 2)$
    \State $\mathbf{k}_4 \leftarrow f(t + h, \mathbf{y} + h \cdot \mathbf{k}_3)$
    \State $\mathbf{y} \leftarrow \mathbf{y} + \frac{h}{6}(\mathbf{k}_1 + 2\mathbf{k}_2 + 2\mathbf{k}_3 + \mathbf{k}_4)$
    \State $t \gets t + h$
    
    \State Добавить $\mathbf{y}$ в $\mathbf{Y}$
\EndFor

\State \Return $\mathbf{Y}$
\end{algorithmic}
\end{algorithm}




Перед обучением траектория нормируется покомпонентно к нулевому
среднему и единичной дисперсии. Затем формируются окна вида
$\bigl(y_t,\;[y_{t+1},\dots,y_{t+L}]\bigr)$ методом скользящего окна.
Полученное множество окон делится в отношении $80/20$ по времени и
последовательности: первые $80\%$ — обучающая выборка, последние
$20\%$ — валидационная. Такой временной, а не случайный, сплит
гарантирует, что модель не «подсматривает» в будущее.

\section{Архитектура модели}

\label{sec:fnode}

В данном разделе описывается архитектура предагаемого подода --- \textbf{FNODE}
(\emph{Fused NODE}). Модель основывается на идеях CSODE, отличаясь малыми архитектурными изменениями 
и \emph{рецептом обучения}. Все полезные приёмы, по отдельности
показавшие положительный эффект на хаотических и квазипериодических
системах, объединены в единый протокол.

\begin{figure}[h]
  \centering
    \includegraphics[width=1\textwidth]{./img/theory/arc.png}
    \caption{Архитектура FNODE}
    \label{fig:csode_p1d}
\end{figure}

Векторное поле модели имеет вид
\[
    f(z) \;=\; -\,\mathrm{softplus}(\gamma) \cdot z \;+\; \mathrm{MLP}(z),
\]
где $\gamma\in\mathbb{R}$ --- обучаемый скалярный параметр сжатия, а
$\mathrm{MLP}$ --- двухслойная SiLU-сеть. Линейный член
$-\mathrm{softplus}(\gamma) \cdot z$ гарантирует диссипативность линейной
части и ограниченность long-horizon-роллаута.

Архитектура FNODE включает в себя:
\begin{itemize}
    \item \textbf{Augmented state}: модель работает в расширенном
    пространстве $z=[y\,\Vert\,0_a]$, как ANODE;

    \item \textbf{Softplus-contraction term}: обучаемый скаляр $\gamma$,
    обеспечивающий мягкую глобальную стягиваемость траекторий
    ($-\mathrm{softplus}(\gamma)<0$ при любом $\gamma$);

    \item \textbf{2-слойный SiLU-MLP}: компактное нелинейное ядро
    векторного поля;

    \item \textbf{Curriculum-горизонт}: длина обучающего окна
    наращивается по расписанию \texttt{"20:1,50:1,85:1,150:2"} ---
    модель сначала учит one-step prediction, затем постепенно
    переходит к мульти-шаговому rollout-у;

    \item \textbf{Teacher-forcing chunks}: rollout длиной $L$
    разбивается на чанки длины $\mathrm{tf\_chunk}=10$, между
    которыми подаётся ground-truth, что снижает накопление ошибки;

    \item \textbf{Spectral FFT loss}: к стандартному MSE добавляется
    штраф на расхождение спектров мощности с весом
    $\lambda_{\mathrm{spec}}=3\cdot 10^{-3}$;


\end{itemize}

Для конфигурации
\[
    aug\_dim=4,
    \qquad
    mlp\_hidd=64,
    \qquad
    \lambda_{\mathrm{spec}}=3\cdot 10^{-3},
\]
число обучаемых параметров FNODE составляет 4\,612 (для системы
с $\dim=3$, итоговый $total\_dim=\dim+aug=7$).

\section{Детали обучения}

\subsection{Конфигурация этапа обучения}

\begin{itemize}
    \item \textbf{Оптимизатор:} Adam~\cite{kingma2017adam} — 
    метод стохастической оптимизации при обучении
    нейронных сетей. Он сочетает идеи накопления первого момента
    градиента\cite{sutskever2013importance} и адаптивного масштабирования
    шага по второму моменту, как в RMSProp;

    \item \textbf{LR Scheduler:} Для контроля learning rate был выбран
    \textit{cosine learning rate scheduler}. Идея данного подхода
    состоит в том, чтобы постепенно уменьшать значение learning rate от
    начального значения до некоторого минимального значения, причём
    уменьшение происходит не линейно, а по косинусоидальному закону;

    \item \textbf{Warmup}: линейный warmup на 50 эпох;
    
    \item \textbf{Gradient clipping}: \texttt{grad\_clip}=1.0
    стабилизирует обратное распространение через RK4-интегратор.

\end{itemize} 

\subsection{Curriculum learning}
\label{sec:curriculum_learning}

При обучении NeuralODE напрямую на длинных траекториях накопленная ошибка быстро уводит
модель с эталонной траектории, это напрямую виляет на качество обучения модели. 
Поэтому в работе используется \emph{curriculum learning}: 
длина обучающего окна увеличивается, что позволяет модели учиться с легких примеров и постепенно переходить к сложным.

На эпохе обучения используется не вся траектория, а её начальный
фрагмент длины
\[
    L = L(\text{epoch}).
\]
В начале обучения берётся короткое окно
\[
    L_{0}=20,
\]
после чего длина окна постепенно возрастает до максимального значения
\[
    L_{\max}\in\{150,200\},
\]
в зависимости от рассматриваемой динамической системы и конкретного
эксперимента. Для системы Хиндмарша-Роуза используется
$L_{\max}=200$, для квазипериодического осциллятора и системы Рёсслера
— $L_{\max}=150$ или $L_{\max}=200$.

Такой подход позволяет сначала обучить модель на коротких и более
стабильных развёртываниях, а затем постепенно усложнять задачу,
увеличивая горизонт прогноза. В результате сеть сначала осваивает
локальную динамику, а затем дообучается на более длинных фрагментах
траектории.

\subsection{Teacher forcing}

Для борьбы с накопленной ошибкой вместе с ~\ref{sec:curriculum_learning} в обучении используется метод \textit{teacher forcing}.

Идея teacher forcing состоит в том, что во время обучения модель периодически получает истинное состояние системы из обучающей выборки вместо собственного предыдущего предсказания. Пусть динамика задаётся как
\[
    \frac{d\mathbf{x}(t)}{dt} = f_\theta(\mathbf{x}(t), t),
\]
где $f_\theta$ --- нейронная сеть. Вместо того чтобы всегда интегрировать траекторию только от предсказанного состояния $\hat{\mathbf{x}}(t_i)$, модель может переинициализироваться истинным значением $\mathbf{x}(t_i)$ и затем предсказывать следующий участок траектории:
\[
    \hat{\mathbf{x}}(t_{i+1}) =
    \mathrm{ODESolve}\bigl(f_\theta, \mathbf{x}(t_i), t_i, t_{i+1}\bigr).
\]

Такой подход стабилизирует обучение, ускоряет сходимость и помогает модели лучше аппроксимировать локальную динамику системы. При этом на этапе инференса модель обычно работает автономно, используя собственные предсказания как начальные условия для дальнейшего интегрирования.

\section{Метрики качества}

Для оценки качества и сравнения исследуемых подходов был сформирован список численных метрик:

\begin{itemize}
    \item \textbf{MSE} — средняя квадратичная ошибка.

    MSE считает среднюю разницу, возведённую в квадрат, между реальным
    и предсказанным значениями в каждом наблюдении. 

    \[
    MSE(y, \hat{y}) = \frac{1}{n} \sum_{i=1}^{n}
    \left( y_i - \hat{y}_i \right)^2.
    \]

    MSE сильно штрафует большие ошибки, что полезно, когда траектории начинают
    сильно расходиться. 
    
    \item \textbf{MAE} — средняя абсолютная ошибка.

    MAE считает среднюю разницу по модулю между реальным и
    предсказанным значениями в каждом наблюдении:

    \[
    MAE(y, \hat{y}) = \frac{1}{n} \sum_{i=1}^{n}
    \mid  y_i - \hat{y}_i  \mid .
    \]

    MAE менее чувствительна к выбросам, что делает её более показательной в
    экспериментах с шумными данными.
 
    \item \textbf{Дивергенция Кульбака–Лейблера}

    Дивергенция Кульбака–Лейблера \cite{kullback1951information} измеряет, насколько сильно предсказанное
    распределение $\hat{y}$ отличается от истинного распределения $y$.
    
    Для дискретных случайных величин с функциями вероятности $p(k)$
    (истинное распределение) и $q(k)$ (предсказанное распределение)
    дивергенция Кульбака–Лейблера определяется как:
    \[
    D_{KL}(P \parallel Q) = \sum_{k} p(k) \, \log \frac{p(k)}{q(k)},
    \]
    где сумма берётся по всем возможным исходам $k$.
    
    При $p(k) \equiv q(k)$ для всех $k$ значение $D_{KL} = 0$ и чем сильнее различаются
    распределения, тем оно больше.

    \item \textbf{Время валидности предсказания}
    $T_{\text{valid}}$ — общее время, когда нормированная RMSE по
    компонентам не превышает порог $\sigma_{\text{th}} = A$.


\end{itemize}

Также для конкретных систем и экспериментов использовались другие специфичные метрики. Их описание будет представлено в разделе с экспериментами.


\chapter{Эксперименты}

Весь код доступен в \url{https://github.com/akon1te/ml-for-dynamic-systems/tree/ode_vkr_epxs}.

\section{Результаты экспериментов}

Модели оценивалась на валидационный частях траекторий каждой системы, выделенных заранее. 

Для систем \emph{Рёсслера} и \emph{Квазипериодического осциллятора} считались метрики, описанные выше. 

Для системы \emph{Хиндмарша-Роуз} вместо $D_{KL}$ считалось время берстов и вычислялась разница $T_{burst\_diff} = T_{burst}\_real - T_{burst}\_pred$. 

\begin{table}[H]
\centering
\resizebox{\textwidth}{!}{%
    \begin{tabular}{c|l|ccc|ccc|ccc}
    \toprule
    \multirow{2}{*}{} &
    \multirow{2}{*}{\textbf{Модели}} &
    \multicolumn{3}{c|}{\textbf{Рёсслер}} &
    \multicolumn{3}{c|}{\textbf{Хиндмарш-Роуз}} &
    \multicolumn{3}{c}{\textbf{КП Осц}} \\
    \cmidrule(lr){3-5}
    \cmidrule(lr){6-8}
    \cmidrule(lr){9-11}
    & &
    MSE & MAE & $D_{\mathrm{KL}} $ &
    MSE & MAE & $T_{burst\_diff} $ &
    MSE & MAE & $D_{\mathrm{KL}} $ \\
    \midrule

    \multirow{2}{*}{A}
    & MLP
      & 8.310 & 7.025 & 0.920
      & & &
      & & & \\
    & RNN
      & 6.606 & 4.904 & 0.781
      & 1.506 & 0.868 & 0.06
      & 2.695 & 1.740 & 0.025 \\

    \midrule

    \multirow{3}{*}{B}
    & NODE
      & \cellcolor{secondblue}3.233 & 2.188 & 0.591
      & \cellcolor{bestgreen}1.279 & \cellcolor{bestgreen}0.717 & 0.07
      & \cellcolor{secondblue}0.235 & \cellcolor{secondblue}0.151 & 0.030 \\
    & ANODE
      & 3.335 & 2.128 & 0.470
      & 1.614 & 0.958 & $-$0.02
      & 0.502 & 0.304 & 0.023 \\
    & CSODE
      & 2.910 & \cellcolor{secondblue}2.012 & \cellcolor{secondblue}0.465
      & 1.559 & 0.907 & \cellcolor{bestgreen}$-$0.014
      & 0.346 & 0.201 & \cellcolor{secondblue}0.020 \\

    \midrule

    & \textbf{FNODE}
      & \cellcolor{bestgreen}2.107 & \cellcolor{bestgreen}1.245 & \cellcolor{bestgreen}0.113
      & \cellcolor{secondblue}1.370 & \cellcolor{secondblue}0.758 & \cellcolor{secondblue}0.016
      & \cellcolor{bestgreen}0.148 & \cellcolor{bestgreen}0.090 & \cellcolor{bestgreen}0.008 \\

    \bottomrule
    \end{tabular}%
}
\caption{Численное сравнение моделей на трех системах}
\end{table}

Модель FNODE при сравнимом количестве параметров показала наилучшие или близкие к лучшим результаты на всех трёх системах. 

На системах \emph{Рёсслера и КП ген} она значительно превосходит остальные методы по точности (MSE, MAE) и расхождению распределений ($D_{\mathrm{KL}}$). 
Для системы \emph{Хиндмарш–Роуз} FNODE незначительно уступает NODE по метрикам MSE/MAE, но демонстрирует более точное воспроизведение бёрстовой динамики по сравнению с другими моделями. 


\newpage
Далее представлены визуализации смоделированных траекторий:

\begin{figure}[H]
    \centering
    \includegraphics[width=0.9\textwidth]{./img/rossler_exp/rossler_phase_full.png}
    \caption{Фазовые портреты на валидационной части траеткории системы Рёсслера}
\end{figure}

\begin{figure}[H]
    \centering
    \includegraphics[width=0.9\textwidth]{./img/hr_exps/hindmarsh_rose_phase_full.png}
    \caption{Фазовые портреты на валидационной части траеткории системы Хиндмарша-Роуза}
\end{figure}

\begin{figure}[H]
    \centering
    \includegraphics[width=0.9\textwidth]{./img/qp_exps/quasi_periodic_phase_val.png}
    \caption{Фазовые портреты на валидационной части траеткории квазипериодического осциллятора}
\end{figure}

\textbf{В главе~\ref{chapter:additional_coparisons} представлены дополнительные сравнения моделей}: распределение предсказаний, графики накапливаемой ошибки от времени моделирования.

\section{Эксперименты с шумом}

В реальных приложениях наблюдаемые временные ряды почти всегда
содержат измерительный шум, поэтому устойчивость модели к шуму является
самостоятельным важным критерием качества. Для хаотических систем эта
проблема особенно существенна: шум в наблюдениях может не только
искажать локальную оценку векторного поля, но и приводить к сильному
расхождению траекторий при авторегрессивном прогнозировании.

В этом разделе исследуется, как точность всех рассматриваемых
архитектур деградирует при добавлении к обучающей траектории
аддитивного шума, и насколько хорошо модели обобщаются на «чистую»
истинную динамику в режиме обучения с зашумлёнными данными.

К нормированной траектории $\tilde y_t \in \mathbb{R}^d$ добавляется
покомпонентный гауссов шум:
\[
    \tilde y_t^{\text{noisy}} = \tilde y_t + \xi_t,
    \qquad \xi_t \sim \mathcal{N}\bigl(0,\,\sigma_n^{2} I_d\bigr),
\]
где $\sigma_n$ задаёт относительный уровень шума. Поскольку
$\tilde y$ имеет единичную дисперсию покомпонентно, $\sigma_n$
соответствует уровню шума относительно амплитуды сигнала.

\vspace{1em}

Качество оценивается \emph{на чистой}, то есть незашумлённой,
валидационной траектории, что позволяет отделить ошибку модели как
таковой от шума в данных. Кроме того, \textbf{основной метрикой будет $T\_valid$} - время когда предсказанная траектория поточечно отличается от таргетной меньше, чем установленный порог.

\begin{table}[H]
\centering
    \begin{tabular}{l|cccc}
    \toprule
    \textbf{Метрика / Модель} &
    \textbf{NODE} &
    \textbf{ANODE} &
    \textbf{CSODE} &
    \textbf{FNODE} \\
    \midrule

    \multicolumn{5}{l}{\textbf{Хиндмарш-Роуз (T\_valid)}} \\
    $\sigma=0.01$ Noise & 40.00 & 8.80 & 51.20 & 40.20 \\
    $\sigma=0.05$ Noise & 39.00 & 8.90 & 39.10 & 38.90 \\
    $\sigma=0.20$ Noise & 9.30 & 6.20 & 9.40 & 9.30 \\

    \midrule

    \multicolumn{5}{l}{\textbf{Рёсслер (T\_valid)}} \\
    $\sigma=0.01$ Noise & 79.25 & 26.40 & 61.65 & 76.70 \\
    $\sigma=0.05$ Noise & 90.80 & 26.40 & 79.30 & 86.20 \\
    $\sigma=0.20$ Noise & 26.35 & 8.90 & 26.40 & 40.35 \\

    \midrule

    \multicolumn{5}{l}{\textbf{КП (T\_valid)}} \\
    $\sigma=0.01$ Noise & 159.98 & 59.42 & 159.98 & 159.98 \\
    $\sigma=0.05$ Noise & 159.98 & 16.32 & 159.98 & 159.98 \\
    $\sigma=0.20$ Noise & 5.58 & 2.08 & 8.00 & 2.58 \\

    \bottomrule
    \end{tabular}
\caption{Численное сравнение моделей на трех системах с добавлением разного уровня шума}
\end{table}

Анализ таблицы показывает, что увеличение уровня шума в обучающих данных приводит к снижению времени корректного прогноза $T_{\mathrm{valid}}$ для всех моделей. Наиболее выраженная деградация наблюдается при $\sigma = 0.20$, что указывает на высокую чувствительность методов к сильному зашумлению данных.

Для системы Хиндмарша--Роуза лучшие значения $T_{\mathrm{valid}}$ в большинстве случаев показывает CSODE, тогда как ANODE существенно уступает остальным методам. 

Для системы Рёсслера наиболее устойчивыми являются NODE и FNODE: NODE даёт лучшие результаты при малом и среднем шуме, а FNODE --- при сильном шуме. 

В случае квазипериодического генератора при малом и среднем шуме NODE, CSODE и FNODE достигают практически максимального времени прогноза, однако с добавлением шума наблюдется сильная деградация способности прогноза моделей.


\section{Эксперименты зависимости качества модели от количества данных}

Дополнительно проводились эксперименты с разным объемом тренировочных данных - на одном и том же $t$ бралась разная часть $(0.05, 0.1, 0.2, 0.5, 1)$ от полной тренировочной выборки.

\begin{figure}[H]
    \centering
    \includegraphics[width=1\textwidth]{./img/rossler_exp/rossler_lowdata.png}
    \caption{$RMSE$ и $D_{KL}$ при разной доле обучающих данных на системе Рёсслера}
\end{figure}

\begin{figure}[H]
    \centering
    \includegraphics[width=1\textwidth]{./img/hr_exps/hr_lowdata.png}
    \caption{$RMSE$ и $D_{KL}$ при разной доле обучающих данных на системе Хиндмарша-Роуза}
\end{figure}

\begin{figure}[H]
    \centering
    \includegraphics[width=1\textwidth]{./img/qp_exps/qp_lowdata.png}
    \caption{$RMSE$ и $D_{KL}$ при разной доле обучающих данных на Квазипериодическом осцилляторе}
\end{figure}


Для системы Рёсслера метод FSODE показывает наиболее устойчивые результаты: значения ошибки остаются сравнительно низкими, а дивергенция $D_{KL}$ уменьшается до минимальных значений при больших объёмах обучающих данных. Метод ANODE, напротив, демонстрирует наибольшие значения как ошибки, так и $D_{KL}$, что указывает на худшее восстановление динамики и распределения состояний.

На системе Хиндмарша--Роуза различия между методами по ошибке менее выражены, однако по метрике $D_{KL}$ преимущество FSODE становится заметным. При средних и больших долях обучающих данных FSODE и CSODE обеспечивают наименьшую дивергенцию, тогда как ANODE сохраняет более высокие значения $D_{KL}$.

Для Квазипериодического осциллятора все методы существенно улучшают качество при увеличении объёма данных. Наилучшее значение ошибки при максимальной доле обучающей выборки достигается методом CSODE, тогда как FSODE и NODE показывают близкие и устойчивые результаты по $D_{KL}$. ANODE во всех экспериментах в среднем уступает остальным подходам.

Таким образом, FSODE можно выделить как наиболее стабильный метод, обеспечивающий хорошее качество восстановления как траекторий, так и распределений состояний. NODE и CSODE также демонстрируют конкурентоспособные результаты, причём CSODE особенно эффективен для квазипериодической динамики. Метод ANODE в рассматриваемых экспериментах не показывает преимуществ и требует дополнительной настройки.


\chapter{Заключение}

В данной работе было проведено исследование возможностей моделей семейства Neural Ordinary Differential Equations для моделирования и прогнозирования нелинейных динамических систем, включая хаотические и квазипериодические режимы. 
Рассматривались классическая NeuralODE, Augmented NeuralODE, ControlSynth NeuralODE, а также предложенная в работе архитектура FNODE, объединяющая аугментацию состояния, несколько MLP-блоков, LayerNorm и SiLU-активацию. 
Тестирование проводилось на трёх системах: системе Рёсслера, модели Хиндмарша–Роуза и Квазипериодическом осцилляторе.

Эксперименты показали, что все рассмотренные NeuralODE-подходы способны аппроксимировать фазовые траектории и воспроизводить статистические свойства аттракторов. 
При этом FNODE продемонстрировала наилучшие результаты по совокупности метрик (MSE, MAE, дивергенция Кульбака–Лейблера) на всех трёх системах, особенно в условиях долгосрочного прогнозирования. 
Использование curriculum learning с постепенным увеличением длины обучающего окна позволило стабилизировать обучение и улучшить качество долгосрочных предсказаний. 
В экспериментах с аддитивным шумом FNODE сохраняла более высокую робастность по сравнению с базовыми архитектурами, что подтверждает её пригодность для задач с зашумлёнными данными.

\subsection*{Ограничения и будущие направления}

Несмотря на достигнутые результаты, работа имеет ряд ограничений. 
Во-первых, исследование ограничено тремя относительно низкоразмерными динамическими системами; поведение моделей на системах большей размерности или с более сложной пространственно-временной структурой требует отдельного изучения. 
Во-вторых, архитектура FNODE, как и другие модели NeuralODE, чувствительна к выбору гиперпараметров, в частности темпа обучения и параметров аугментации, что может потребовать тщательного подбора для новых задач. 
В-третьих, вычислительная стоимость обучения остаётся значительной из-за необходимости многократного численного интегрирования.


\putbibliography


%====================================================================


\chapter{Дополнительные репрезентативные сравнения подходов}
\label{chapter:additional_coparisons}

\subsection*{Система Рёсслера}

\begin{figure}[H]
    \centering
    \includegraphics[width=0.75\textwidth]{img/rossler_exp/rossler_hist_comp0.png}
    \vspace{1pt} 
    \includegraphics[width=0.75\textwidth]{img/rossler_exp/rossler_hist_comp1.png}
    \caption{Сопоставление распределений предсказаний моделей с истинным распределением целевых переменных x, y на тестовой выборке.}
\end{figure}

\begin{figure}[H]
    \centering
    \includegraphics[width=0.9\textwidth]{img/rossler_exp/rossler_pointwise_err.png}
    \caption{Накапливаемая ошибка предсказаний моделей.}
\end{figure}



\subsection*{Система Хиндмарша-Роуза}

\begin{figure}[H]
    \centering
    \includegraphics[width=0.9\textwidth]{img/hr_exps/hindmarsh_rose_hist_comp0.png}
    \vspace{1pt} 
    \includegraphics[width=0.9\textwidth]{img/hr_exps/hindmarsh_rose_hist_comp2.png}
    \caption{Сопоставление распределений предсказаний моделей с истинным распределением целевых переменных x, z на тестовой выборке.}
\end{figure}

\begin{figure}[H]
    \centering
    \includegraphics[width=0.9\textwidth]{img/hr_exps/hindmarsh_rose_pointwise_err.png}
    \caption{Накапливаемая ошибка предсказаний моделей.}
\end{figure}



\subsection*{Квазипериодический осциллятор}

\begin{figure}[H]
    \centering
    \includegraphics[width=0.9\textwidth]{img/qp_exps/quasi_periodic_hist_comp2.png}
    \caption{Сопоставление распределений предсказаний моделей с истинным распределением целевых переменных Z на тестовой выборке.}
\end{figure}

\begin{figure}[H]
    \centering
    \includegraphics[width=0.9\textwidth]{img/qp_exps/quasi_periodic_pointwise_err.png}
    \caption{Накапливаемая ошибка предсказаний моделей.}
\end{figure}


\section*{Дополнительные метрики}


\begin{table}[H]
    \centering
    \label{tab:spectral_l1}
    \begin{tabular}{lccc}
        \toprule
        Модель & Рёсслер & Хиндмарш-Роуз & КГ Осц \\
        \midrule
        RNN   & 2.4361 & 2.0557 & 2.2680 \\
        NODE  & \cellcolor{bestgreen}0.6877 & 3.8823 & \cellcolor{secondblue}0.1140 \\
        ANODE & 1.1208 & \cellcolor{secondblue}1.6216 & 0.1881 \\
        CSODE & 1.1208 & 3.1523 & 0.1881 \\
        FNODE & \cellcolor{secondblue}0.8298 & \cellcolor{bestgreen}1.2928 & \cellcolor{bestgreen}0.0439 \\
        \bottomrule
    \end{tabular}
    \caption{Спектральная ошибка L1 ($\downarrow$) в динамических системах и моделях.
              Более низкие значения указывают на лучшее спектральное соответствие между истинной и прогнозируемой траекториями.}
\end{table}


\begin{table}[H]
    \centering
    \label{tab:w1_rossler}
    \begin{tabular}{lc}
        \toprule
        Модель & $W_1$ \\
        \midrule
        RNN   & 0.1245\,/\,0.1496\,/\,0.1241 \\
        NODE  & 0.1607\,/\,0.1654\,/\,0.1416 \\
        ANODE & \cellcolor{secondblue}0.1217\,/\,0.1269\,/\,0.1084 \\
        CSODE & \cellcolor{secondblue}0.1217\,/\,0.1269\,/\,0.1084 \\
        FNODE & \cellcolor{bestgreen}0.0251\,/\,0.0171\,/\,0.0276 \\
        \bottomrule
    \end{tabular}
    \caption{Расстояние Вассерштейна-1 ($\downarrow$) для системы Рёсслера.
              Значения представлены в виде $W_1(x_1)\,/\, W_1(x_2)\,/\, W_1(x_3)$.}
    
\end{table}

\begin{table}[H]
    \centering
    \label{tab:w1_hr}
    \begin{tabular}{lc}
        \toprule
        Модель & $W_1$ \\
        \midrule
        RNN   & 0.0552\,/\,0.1078\,/\,0.2690 \\
        NODE  & \cellcolor{secondblue}0.0778\,/\,0.1548\,/\,0.0979 \\
        ANODE & 0.0423\,/\,0.1185\,/\,0.1986 \\
        CSODE & \cellcolor{bestgreen}0.0509\,/\,0.0605\,/\,0.1489 \\
        FNODE & 0.0676\,/\,0.1348\,/\,0.2805 \\
        \bottomrule
    \end{tabular}
    \caption{Расстояние Вассерштейна-1 ($\downarrow$) для системы Хиндмарша--Роуза.
              Значения представлены в виде $W_1(x_1)\,/\, W_1(x_2)\,/\, W_1(x_3)$.}
    
\end{table}


\begin{table}[H]
    \centering
    \label{tab:w1_qp}
    \begin{tabular}{lc}
        \toprule
        Модель & $W_1$ \\
        \midrule
        RNN   & 0.0303\,/\,0.0297\,/\,0.0681 \\
        NODE  & 0.0062\,/\,0.0063\,/\,0.0149 \\
        ANODE & \cellcolor{secondblue}0.0056\,/\,0.0051\,/\,0.0133 \\
        CSODE & \cellcolor{secondblue}0.0056\,/\,0.0051\,/\,0.0133 \\
        FNODE & \cellcolor{bestgreen}0.0048\,/\,0.0050\,/\,0.0049 \\
        \bottomrule
    \end{tabular}
    \caption{Расстояние Вассерштейна-1 ($\downarrow$) для системы квазипериодического осциллятора.
              Значения представлены в виде $W_1(x_1)\,/\, W_1(x_2)\,/\, W_1(x_3)$.}
    
\end{table}


\begin{table}[H]
\centering
    \setlength{\tabcolsep}{4pt}
    \begin{tabular}{l|c|c|c|c|}
    \toprule
    \textbf{Критерий} & \textbf{NODE} & \textbf{SINDy} & \textbf{Seq2Seq} & \textbf{ESN} \\
    \midrule
    Непрерывность времени & \checkmark & -- & -- & -- \\
    Неравномерные данные & \checkmark & -- & -- & -- \\
    Долгосрочный прогноз & \checkmark & -- & \checkmark  & \checkmark \\
    Интерпретируемость & $\triangle$ & \checkmark & -- & -- \\
    Устойчивость к шуму & $\triangle$ & -- & $\triangle$ & \checkmark \\
    \midrule
    \multicolumn{5}{l}{\footnotesize \checkmark — хорошо, $\triangle$ — средне, -- — плохо/нет} \\
    \bottomrule
    \end{tabular}%
\end{table}


\chapter*{Использование больших языковых моделей}

При подготовке данной выпускной квалификационной работы большие языковые модели (LLM), в частности Claude\footnote{\url{https://claude.com/}} и Deepseek\footnote{\url{https://www.deepseek.com/}}, использовались как инструменты редактирования и повышения качества оформления работы. Их применение включало:

\begin{itemize}
    \item \textbf{Генерацию шаблонов LaTeX и помощь с форматированием таблиц и разделов документов}, например создание или уточнение структуры разделов, синтаксиса размещения рисунков и настройки окружений;

    \item \textbf{Проверку грамматики и пунктуации}, обеспечивающую незначительную редакторскую поддержку для повышения ясности текста и соблюдения академического стиля.
\end{itemize}

Большие языковые модели не использовались для формирования исследовательских идей, интерпретации результатов или написания содержательной части работы. Все ключевые научные решения и выводы принадлежат автору.


\end{document}
```
